\documentclass[11pt,leqno]{amsart}

\usepackage[T1]{fontenc}
\usepackage[utf8]{inputenc}
\usepackage{lmodern}
\usepackage{eulervm}
\usepackage{amsmath,amssymb,amsthm}

\usepackage{booktabs,array}
\usepackage{enumitem}
\usepackage{xurl}
\usepackage[colorlinks,linkcolor=blue,citecolor=blue]{hyperref}
\usepackage[nameinlink,noabbrev]{cleveref}
\usepackage{orcidlink}
\setlength{\emergencystretch}{2em}


\newtheorem{theorem}{Theorem}[section]
\newtheorem{proposition}[theorem]{Proposition}
\newtheorem{lemma}[theorem]{Lemma}
\newtheorem{corollary}[theorem]{Corollary}
\newtheorem{conjecture}[theorem]{Conjecture}
\theoremstyle{remark}
\newtheorem{remark}[theorem]{Remark}

\numberwithin{equation}{section}


\newcommand{\Q}{\mathbb Q}
\newcommand{\Z}{\mathbb Z}
\newcommand{\F}{\mathbb F}
\newcommand{\OK}{\mathcal O_K}
\DeclareMathOperator{\Cl}{Cl}
\DeclareMathOperator{\Nm}{N}
\DeclareMathOperator{\ord}{ord}
\DeclareMathOperator{\rank}{rank}

% The journal's 2006 class predates the 2020 MSC edition.
\makeatletter
\@namedef{subjclassname@2020}{2020 Mathematics Subject Classification}
\makeatother

\title[Prescribed ideal classes and elliptic points]{Prescribed ideal classes and elliptic points in Ennola cubic fields}
\author[J.~Lu]{Junyu Lu\,\orcidlink{0009-0001-4973-3198}}
\address{Department of Mathematics, Yang-En University, Majia Town, Luojiang District, Quanzhou City, Fujian Province, China, 362014}
\email{jy.lu@outlook.com}
\date{}
\subjclass[2020]{Primary 11R29; Secondary 11G05, 11R16, 11D09}
\keywords{Ennola cubic field, ideal class group, elliptic curve, squareclass, Pell equation}

\begin{document}
\begin{abstract}
  For every prime $\ell\geq5$ and integer $n\geq2$, we construct an Ennola cubic field with a specified prime above $\ell$ of exact ideal-class order $2^n$. The same squareclass obstructions give four independent rational points on associated elliptic curves. A Pell equation describes all parameters satisfying the required square and congruence conditions. A genus-one base change gives five independent sections whose generated subgroup is saturated at two. It also yields infinitely many geometrically nonisomorphic rational fibres of rank at least five.
\end{abstract}
\maketitle


\section{Introduction}

Ennola's family of cubic polynomials arises in the study of exceptional units. Recall that a unit $\varepsilon$ is exceptional if $\varepsilon-1$ is also a unit \cite{Ennola1991,Louboutin2017}. For an integer $a\geq3$, let $\theta$ be the largest real root of
\[
  f_a(X)=X^3+(a-1)X^2-aX-1=X(X-1)(X+a)-1
\]
and put $K_a=\Q(\theta)$. This choice fixes the real embedding in which both $\theta$ and $\theta-1$ are positive. The polynomial is irreducible, has three real roots, and has positive nonsquare discriminant $(a^2+3a-1)^2-32$. Thus $K_a$ is a totally real, non-Galois cubic field. The identity $\theta(\theta-1)(\theta+a)=1$ exhibits $\theta$ as an exceptional unit; both $\theta$ and $\theta-1$ have norm one. The elementary field properties are recalled in Section~2.

Ennola conjectured that $\theta$ and $\theta-1$ form a fundamental system of units of the maximal order; see \cite{Ennola1991,Louboutin2017,ChoiKim2025}. We require only the known description of their unit squareclasses, so our results do not depend on this conjecture. We write $\Cl(K_a)$ for the ordinary ideal class group of the maximal order $\mathcal O_{K_a}$. The same squareclasses control principal ideal powers and the independence of rational points on
\begin{equation}\label{eq:elliptic}
  E_a:\quad y^2=x(x+1)(x-a)+1.
\end{equation}
The connection is the identity
\begin{equation}\label{eq:common-norm}
  \Nm_{K_a/\Q}(x+\theta)=x(x+1)(x-a)+1.
\end{equation}
For a number field $L$, $\Nm_{L/\Q}$ denotes the field norm. For a nonzero integral ideal $\mathfrak c\subset\mathcal O_L$, set $\Nm\mathfrak c=\#(\mathcal O_L/\mathfrak c)$, and extend this norm multiplicatively to nonzero fractional ideals. Brackets denote element squareclasses in $K_a^\times/K_a^{\times2}$ or ideal classes in $\Cl(K_a)$, according to the argument. Finite-place valuations are normalized to have value one on a uniformizer. The point at infinity is the identity $O$ of $E_a$.

Our first result prescribes both the rational prime and the exact order of a prime ideal class.

\begin{theorem}\label{thm:prescribed-prime}
  Let $\ell\geq5$ be a rational prime. For integers $n\geq2$, put $N=2^n$ and
  \[
    (D,a)=
    \begin{cases}
      \left(4,(\ell^N-73)/12\right),     & \ell\ne137, \\[2pt]
      \left(24,(137^N-14353)/552\right), & \ell=137.
    \end{cases}
  \]
  Then $a$ is an integer with $a\geq3$ and $a\equiv2\pmod4$. Moreover, the ideal
  \[
    \mathfrak p=(\ell,D\theta-D+1)\mathcal O_{K_a}
  \]
  is prime of norm $\ell$, and
  \[
    \mathfrak p^N=(D\theta-D+1),\qquad
    \ord_{\Cl(K_a)}[\mathfrak p]=N.
  \]
  For each fixed $\ell$, these parameters give infinitely many isomorphism classes of fields $K_a$ as $n$ varies.
\end{theorem}

The change of linear form at $137$ is needed for the stated prime-ideal construction; see Remark~\ref{rem:137}. The defining order need not be maximal. Proposition~\ref{prop:nonmaximal} gives an explicit infinite sequence where it is not.

\begin{theorem}\label{thm:four-points}
  Suppose that $a\geq3$, $a\equiv30$ or $34\pmod{52}$. If there are positive integers $u,v$ satisfying
  \begin{equation}\label{eq:two-squares}
    u^2=12a+73,\qquad v^2=240a+4321,
  \end{equation}
  then the points
  \[
    P=(0,1),\quad Q=(-1,1),\quad
    S=(-3/4,u/8),\quad T=(-15/16,v/64)
  \]
  are independent modulo torsion in $E_a(\Q)$. The subgroup $H$ generated by $P,Q,S,T$ is saturated at two in the sense that $2R\in H$ for $R\in E_a(\Q)$ implies $R\in H$.

  If also $137\nmid u$ and $\gcd(v,23\cdot2503)=1$, then
  \[
    \mathfrak a=(u,4\theta-3)\mathcal O_{K_a},\qquad
    \mathfrak b=(v,16\theta-15)\mathcal O_{K_a}
  \]
  are ideals of the maximal order with
  \[
    \Nm\mathfrak a=u,\quad \mathfrak a^2=(4\theta-3),\qquad
    \Nm\mathfrak b=v,\quad \mathfrak b^2=(16\theta-15),
  \]
  and their classes generate a subgroup isomorphic to $(\Z/2\Z)^2$ in $\Cl(K_a)$.
\end{theorem}

The square conditions define $v^2-20u^2=2861$. Theorem~\ref{thm:pell} describes all positive solutions giving either required congruence on $a$, and Proposition~\ref{prop:subsequence} provides infinitely many parameters satisfying the ideal-class hypotheses. The smallest parameter is $a=134$, with $(u,v)=(41,191)$. We also obtain an explicit family over $\Q(t)$ of rank at least four in Corollary~\ref{cor:rational-family}. These theorems give rank lower bounds and saturation at two; they do not determine ranks uniformly or give full Mordell--Weil bases of the specializations.

Proposition~\ref{prop:rank-five-rational} gives infinitely many rational parameters with rank at least five by a further genus-one base change. This construction has only one integral value of $a$, as shown in Remark~\ref{rem:rank-five-integrality}.

\subsection*{Relation to earlier work}

Ennola proved the unit-squareclass result and constructed $K_a(\sqrt{\theta(\theta-1)})/K_a$ as an unramified quadratic extension for $a\equiv2\pmod4$ \cite[Propositions 3.3 and 6.1]{Ennola1991}; his Theorem~6.1 also gives families whose class groups have two-rank at least two for all sufficiently large admissible parameters. We recall Louboutin's short proof of the squareclass result \cite[Lemma 6]{Louboutin2017} and obtain a second independent unramified radical in Corollary~\ref{cor:hilbert}. The use of signatures and dyadic congruences follows the classical strategy of G.~Gras and M.-N.~Gras \cite[\S I.6]{GrasGras1975}, with the local conditions checked directly here.

Lee \cite[Theorem 1.1]{Lee2010} constructs ideal classes of any prescribed order in infinitely many fields of the same cubic family, written using $\rho=1/(\theta-1)$ and parameter $a+1$. Theorem~\ref{thm:prescribed-prime} fixes the rational prime in advance and gives uniform formulas for a prime above it of exact power-of-two class order. The norm bound of Kala--Sgallov\'a--Tinkov\'a \cite[\S3.5, Proposition 3.7 in arXiv version 2]{KalaSgallovaTinkova2025} gives an alternative exact-order argument for the first parameter formula when $\mathcal O_{K_a}=\Z[\theta]$. Under the same assumption, Louboutin's earlier bound \cite[Theorem 10]{Louboutin2001} gives this argument for $a\geq59$. Our local proof also covers the nonmaximal orders in Proposition~\ref{prop:nonmaximal}. Related power-residue constructions occur in Pincus--Washington \cite[\S4]{PincusWashington2023}.

Connections between elliptic points and cubic ideal classes appear in Duquesne \cite{Duquesne2001} and Laoudi \cite{Laoudi2006}; Buell--Ennola \cite{BuellEnnola1995} study three-primary classes in the quadratic resolvent of this family. For related elliptic curves, Walsh gives a Pell construction of rank at least three for sufficiently large parameters under auxiliary irreducibility hypotheses \cite[Theorem 1.1]{Walsh2023}, and a four-point construction with computational evidence of independence \cite[\S2, Type III]{Walsh2024}. Here we prove simultaneous rank-four and cubic ideal-class results under explicit congruence conditions, and classify all parameters in the resulting Pell construction.

\section{Unit squareclasses and local obstructions}

The rational-root test proves that $f_a(X)$ is irreducible, since $f_a(1)=-1$ and $f_a(-1)=2a-3$. Its three roots lie respectively in $(-a,1-a)$, $(-1,0)$, and $(1,2)$, as follows by evaluation at the six endpoints. In this order of the real embeddings, both $\theta$ and $\theta-1$ have signature $(-,-,+)$ and norm one. Write
\begin{equation}\label{eq:discriminant}
  d_a=\operatorname{disc}(f_a)=(a^2+3a-1)^2-32.
\end{equation}
This integer is positive and odd. Since $\operatorname{disc}(K_a)=d_a/[\mathcal O_{K_a}:\Z[\theta]]^2$, the index is odd and $\mathcal O_{K_a}\otimes\Z_2=\Z_2[\theta]$. The displayed Weierstrass equation has discriminant $16d_a$, so it is nonsingular. Also $d_a$ lies strictly between the consecutive squares $(a^2+3a-2)^2$ and $(a^2+3a-1)^2$. The field discriminant has the same nonsquare squareclass, so the cubic field is non-Galois. See \cite{Ennola1991} for more details.

\begin{lemma}\label{lem:units}
  For every $a\geq3$, the positive-norm units of $K_a$, modulo squares of units, have the four distinct representatives
  \[
    1,\quad \theta,\quad \theta-1,\quad \theta(\theta-1).
  \]
\end{lemma}

\begin{proof}
  We recall the argument in \cite[Lemma 6]{Louboutin2017}. The signs exclude $\theta$ and $\theta-1$ from being squares. Their product has the same squareclass as
  \[
    \epsilon=\frac{\theta}{\theta-1}=\theta^2+a\theta+1.
  \]
  Since $\theta=\epsilon/(\epsilon-1)$, this unit generates $K_a$; its minimal polynomial is $Z^3-(a+4)Z^2+(a+3)Z-1$. If $\epsilon=\eta^2$ in $K_a$, then $\eta$ is an integral unit of degree three. Changing its sign gives norm one, so its minimal polynomial has the form $Z^3-sZ^2+wZ-1$ with $s,w\in\Z$. Comparing the elementary symmetric functions of its squared conjugates gives
  \[
    s^2-2w=a+4,\qquad w^2-2s=a+3,
    \qquad (s-w)(s+w+2)=1.
  \]
  The integer pairs are $(s,w)=(0,-1),(-2,-1)$, giving $a=-2,2$, respectively. Neither is allowed. Thus $[\theta]$ and $[\theta-1]$ are independent in $K_a^\times/K_a^{\times2}$. A square root of a unit is a unit, by its valuations at finite primes. Dirichlet's unit theorem gives eight unit squareclasses, and $\Nm_{K_a/\Q}$ maps this group onto $\{1,-1\}$ because $\Nm_{K_a/\Q}(-1)=-1$. Its kernel therefore consists of the four displayed classes.
\end{proof}

\begin{lemma}\label{lem:obstruction}
  Suppose that $a\geq3$ and $a\equiv2\pmod4$. If $D$ is a positive integer divisible by four, then $D\theta-D+1$ is not a unit times a square in $K_a$.
\end{lemma}

\begin{proof}
  The odd index gives
  \[
    \mathcal O_{K_a}/4\mathcal O_{K_a}\simeq
    (\Z/4\Z)[X]/(X^3+X^2+2X-1).
  \]
  Its reduction modulo two is $\F_8$. The unique square roots of $\theta$ and $\theta-1$ in $\F_8$ are $\theta^2+\theta+1$ and $\theta^2+\theta$. Lifts of the same residue modulo two have equal squares modulo four, whereas
  \begin{align*}
    (\theta^2+\theta+1)^2 & \equiv\theta+2\not\equiv\theta\pmod4,              \\
    (\theta^2+\theta)^2   & \equiv2\theta^2+3\theta+1\not\equiv\theta-1\pmod4.
  \end{align*}
  The second difference is $2(\theta^2+\theta+1)\ne0$ in the free basis. Thus neither $\theta$ nor $\theta-1$ is a square modulo four in the maximal order.

  Put $\gamma=D\theta-D+1$. Since $D\geq4$, the root intervals give its signature as $(-,-,+)$. Thus $\Nm_{K_a/\Q}(\gamma)>0$, and $\gamma\equiv1\pmod{4\mathcal O_{K_a}}$. The products $\gamma$ and $\gamma\theta(\theta-1)$ have negative conjugates. The products $\gamma\theta$ and $\gamma(\theta-1)$ are nonsquares modulo four. These tests exclude squares in $K_a$, since a square root of an algebraic integer is integral. If $\gamma=\varepsilon\xi^2$ for a unit $\varepsilon$, taking norms to $\Q$ gives $\Nm_{K_a/\Q}(\varepsilon)=1$. Lemma~\ref{lem:units} would then make one of the four excluded products a square.
\end{proof}

Put $\alpha=4\theta-3$ and $\beta=16\theta-15$. Their norms are
\begin{equation}\label{eq:alpha-beta}
  \Nm_{K_a/\Q}(\alpha)=12a+73,\qquad \Nm_{K_a/\Q}(\beta)=240a+4321.
\end{equation}

\begin{lemma}\label{lem:four-classes}
  If $a\geq3$ and $a\equiv30$ or $34\pmod{52}$, the four classes $[\theta]$, $[\theta-1]$, $[\alpha]$, and $[\beta]$ are independent in $K_a^\times/K_a^{\times2}$.
\end{lemma}

\begin{proof}
  Both congruences give $a\equiv2\pmod4$. Lemma~\ref{lem:units} excludes the three nontrivial products of the unit classes, and Lemma~\ref{lem:obstruction} excludes the eight products involving exactly one of $\alpha,\beta$. The products $\alpha\beta\theta$ and $\alpha\beta(\theta-1)$ have negative conjugates. It remains to exclude $\alpha\beta$ and $\alpha\beta\theta(\theta-1)$.

  If $a\equiv30\pmod{52}$, then $a\equiv4\pmod{13}$ and
  \[
    f_a(X)\equiv(X-6)(X^2+9X+11)\pmod{13},\qquad d_a\equiv8\pmod{13}.
  \]
  The prime $\mathfrak l=(13,\theta-6)$ of the maximal order has residue field $\F_{13}$. At this prime,
  \[
    \alpha\beta\equiv11,\qquad \theta(\theta-1)\equiv4.
  \]
  The two remaining products reduce to $11$ and $5$, both nonsquares in $\F_{13}$.

  If $a\equiv34\pmod{52}$, then $a\equiv8\pmod{13}$ and
  \[
    f_a(X)\equiv(X-3)(X-4)(X-12)\pmod{13}.
  \]
  These distinct roots define degree-one primes of the maximal order. At the prime with $\theta\equiv3$, we have $\alpha\beta\equiv11$. At the prime with $\theta\equiv12$, we have $\alpha\beta\theta(\theta-1)\equiv5$. Both residues are nonsquares. Thus all fifteen nontrivial products are excluded in either case.
\end{proof}

\section{Prime ideals of prescribed order}

We first record the local step that allows us to work in the maximal order without assuming $\mathcal O_K=\Z[\theta]$.

\begin{lemma}\label{lem:local-power}
  Let $f\in\Z[X]$ be monic and irreducible, let $K=\Q(\theta)$ with $f(\theta)=0$, and let $D,c,m,N$ be integers with $D,m>0$, $N\geq1$, and $\gcd(D,m)=1$. If $\Nm_{K/\Q}(D\theta-c)=m^N$ and $f'(c/D)$ is a $p$-adic unit for every prime $p\mid m$, then
  \[
    I=(m,D\theta-c)\mathcal O_K
    \quad\text{satisfies}\quad \Nm I=m,\qquad I^N=(D\theta-c).
  \]
  If $m$ is a rational prime, $I$ is prime.
\end{lemma}

\begin{proof}
  Fix $p\mid m$. The norm identity implies $f(c/D)\equiv0\pmod p$. The derivative hypothesis and Hensel's lemma give $f(X)=(X-r)h(X)$ in $\Z_p[X]$, with $r\equiv c/D\pmod p$ and $h(r)\in\Z_p^\times$. The coprime factors give $\Z_p[X]/(f)\simeq\Z_p\times\Z_p[X]/(h)$. Taking integral closures in $K\otimes_{\Q}\Q_p$ preserves this product. The first factor is already integrally closed and equals $\Z_p$, so its prime $\mathfrak p_p$ has ramification index and residue degree both one. In the complementary factor, $DX-c$ is a unit modulo $p$ because $h(c/D)$ is a $p$-adic unit. Its image is therefore a unit in the integral closure. Thus $\mathfrak p_p$ is the only prime above $p$ that divides $(D\theta-c)$.

  No prime over a rational prime outside $m$ divides $(D\theta-c)$, by its norm. Comparing valuations of the norm therefore gives
  \[
    (D\theta-c)=\prod_{p\mid m}\mathfrak p_p^{Nv_p(m)}.
  \]
  Taking the ideal sum with $(m)$ replaces each exponent by $v_p(m)$. This proves both identities and the last assertion.
\end{proof}

For later use, direct evaluation gives
\begin{equation}\label{eq:general-norm}
  \Nm_{K_a/\Q}(D\theta-D+1)=D(D-1)a+D^3+D^2-2D+1.
\end{equation}

\begin{proof}[Proof of Theorem~\ref{thm:prescribed-prime}]
  First let $\ell\ne137$ and $D=4$. Because $N$ is divisible by four and $\ell$ is coprime to six, $\ell^N\equiv1\pmod{48}$. Thus $a=(\ell^N-73)/12$ is an integer with $a\geq46$ and $a\equiv2\pmod4$. Equations~\eqref{eq:alpha-beta} and
  \[
    f_a'(3/4)=\frac{2\ell^N-137}{48}
  \]
  give the hypotheses of Lemma~\ref{lem:local-power}. Hence the indicated ideal $\mathfrak p$ has norm $\ell$ and satisfies $\mathfrak p^N=(4\theta-3)$.

  For $\ell=137$ and $D=24$, the congruence $137^4\equiv1\pmod{2208}$ gives
  \[
    a=(137^N-14353)/552\in\Z,\qquad a\equiv2\pmod4,
  \]
  and $a\geq638154$. Equation~\eqref{eq:general-norm} and differentiation give
  \[
    \Nm_{K_a/\Q}(24\theta-23)=137^N,\qquad
    f_a'(23/24)=\frac{22\cdot137^N-304657}{13248}.
  \]
  The denominator is prime to $137$, and $304657\equiv106\pmod{137}$. Lemma~\ref{lem:local-power} again gives a prime of norm $137$ with the required power identity.

  In either case, $[\mathfrak p]^N=1$. If its order were a proper divisor of the power of two $N$, then $\mathfrak p^{N/2}$ would be principal. Writing $\mathfrak p^{N/2}=(\xi)$ would make $D\theta-D+1$ a unit times $\xi^2$, contrary to Lemma~\ref{lem:obstruction}. Hence $[\mathfrak p]$ has exact order $N$. Finally, a finite set of number fields has class groups of bounded exponent. The unbounded orders $2^n$ prove the assertion about field isomorphism classes.
\end{proof}

For comparison, suppose that $\ell\ne137$ and $\mathcal O_{K_a}=\Z[\theta]$. If $\mathfrak p^{N/2}=(\gamma)$, then $\gamma$ is integral and
\[
  |\Nm_{K_a/\Q}(\gamma)|=\ell^{N/2}=\sqrt{12a+73}<2a-3,
\]
since $a\geq46$ and $(2a-3)^2-(12a+73)=4(a-8)(a+2)>0$. The norm bound in \cite[\S3.5, Proposition 3.7 in arXiv version 2]{KalaSgallovaTinkova2025} forces $\gamma$ to be a unit times a rational integer. Its absolute norm would then be a cube, contrary to $3\nmid N/2=2^{n-1}$. This gives the alternative exact-order argument mentioned in the introduction.

Louboutin's bound \cite[Theorem 10]{Louboutin2001} gives the same deduction for $a\geq59$, with his parameter $m=a+1$ and bound $2m-5=2a-3$. The identities $\rho=1/(\theta-1)=\theta(\theta+a)$ and $\theta=\rho^2-(a+1)\rho-(a+1)$ give $\Z[\rho]=\Z[\theta]$, so the maximality hypothesis agrees in the two sets of coordinates.

\begin{corollary}\label{cor:composite}
  Let $n\geq2$, $N=2^n$, and $m\geq5$ with $\gcd(m,6\cdot137)=1$. For $a=(m^N-73)/12$, the ideal $I=(m,4\theta-3)\mathcal O_{K_a}$ has norm $m$, satisfies $I^N=(4\theta-3)$, and has exact class order $N$.
\end{corollary}

\begin{proof}
  The fourth-power congruence modulo $48$ and the derivative calculation in the first part of the preceding proof hold for this $m$ at every prime divisor. Apply Lemmas~\ref{lem:local-power} and \ref{lem:obstruction}, and the same exact-order argument.
\end{proof}

\begin{remark}\label{rem:137}
  The first formula cannot be extended to $\ell=137$ with the same ideal. Put $q=137^N$, $a=(q-73)/12$, and $\alpha=4\theta-3$. The minimal polynomial of $\alpha$ is
  \[
    Z^3+\frac{q-58}{3}Z^2+\frac{2q-137}{3}Z-q.
  \]
  Write this polynomial as $G(Z)$ and put $p=137$. The three normalized polynomials have reductions
  \begin{align*}
    G(p^{N-1}Z)/p^N &\equiv -Z/3-1, \\
    G(pZ)/p^2 &\equiv -Z(58Z+1)/3, \\
    G(Z) &\equiv Z^2(Z-58/3) \pmod p.
  \end{align*}
  Their respective nonzero roots $-3$, $-1/58$, and $58/3$ are simple.

  Hensel's lemma gives three roots of $G$ in $\Q_{137}$ of valuations $N-1$, $1$, and $0$. Thus $137$ splits completely, $(\alpha)=\mathfrak p^{N-1}\mathfrak q$ for two distinct primes of norm $137$, and $(137,\alpha)$ has norm $137^2$.

  For the second parameter formula $a=(q-14353)/552$ in Theorem~\ref{thm:prescribed-prime}, the form $24\theta-23$ avoids this obstruction. At $x=-23/24$, its norm identity gives $y^2=137^N/24^3$, which is not a rational square. Thus it does not give a rational point on \eqref{eq:elliptic} by the denominator-clearing construction used for $4\theta-3$.
\end{remark}

\begin{proposition}\label{prop:nonmaximal}
  For $n\geq2$ and $a_n=(5^{2^n}-73)/12$, if $n\equiv13\pmod{110}$, then
  \[
    23\mid[\mathcal O_{K_{a_n}}:\Z[\theta]].
  \]
  These parameters include infinitely many isomorphism classes of fields.
\end{proposition}

\begin{proof}
  The modular identities
  \[
    5^{506}\equiv1\pmod{6348},\quad
    2^{110}\equiv1\pmod{253},\quad
    5^{2^{13}}\equiv2845\pmod{6348}
  \]
  give $a_n\equiv231\pmod{529}$. Indeed, $2^n\equiv2^{13}\pmod{253}$, and both exponents are even, giving the congruence modulo $506$. Write $a_n=231+529k$. Substitution using $f_{a_n}(\theta)=0$ shows that $\xi=(\theta^2-13\theta+30)/23$ satisfies
  \begin{align*}
    \xi^3 & -(12167k^2+10925k+2454)\xi^2                    \\
          & +(25392k^2+22802k+5119)\xi-(138k+61)(90k+41)=0.
  \end{align*}
  This monic equation makes $\xi$ integral, while the basis $1,\theta,\theta^2$ shows that $\xi\notin\Z[\theta]$ and $23\xi\in\Z[\theta]$. This proves the index assertion. Theorem~\ref{thm:prescribed-prime} gives unbounded class orders along the progression, so finitely many fields cannot account for it.
\end{proof}

\section{Four points and two ideal classes}

\begin{lemma}\label{lem:virtual-units}
  Let $K=K_a$, with $a\geq3$, and let
  \[
    E_K^{(1)}=\ker(\Nm_{K/\Q}:\mathcal O_K^\times\to\{\pm1\})
  \]
  be its group of norm-one units. Define
  \[
    V_K=\{\gamma\in K^\times:\Nm_{K/\Q}(\gamma)>0,\ v_{\mathfrak p}(\gamma)\in2\Z\ \text{for every finite prime }\mathfrak p\}/K^{\times2}.
  \]
  Then there is an exact sequence
  \[
    1\longrightarrow E_K^{(1)}/(\mathcal O_K^\times)^2
    \longrightarrow V_K\longrightarrow\Cl(K)[2]\longrightarrow1.
  \]
  The unit kernel has dimension two over $\F_2$, with basis $[\theta],[\theta-1]$.
\end{lemma}

\begin{proof}
  The even valuations determine a unique fractional ideal $\mathfrak c$ with $(\gamma)=\mathfrak c^2$. Replacing $\gamma$ by a square multiple changes $\mathfrak c$ by a principal ideal, so the map is well defined. Its kernel consists of the classes of norm-one units: if $\mathfrak c=(\xi)$, then $\gamma/\xi^2$ is such a unit. Conversely, each norm-one unit represents an element of the kernel. A square root in $K$ of a unit is a unit, so the left map is injective. If $[\mathfrak c]\in\Cl(K)[2]$, write $\mathfrak c^2=(\gamma)$; replacing $\gamma$ by $-\gamma$ when necessary makes its norm positive because $[K:\Q]=3$. This proves surjectivity. The kernel assertion follows from Lemma~\ref{lem:units}.
\end{proof}

\begin{proof}[Proof of Theorem~\ref{thm:four-points}]
  Define
  \[
    \delta:E_a(\Q)\longrightarrow K_a^\times/K_a^{\times2},
    \qquad\delta(O)=1,\quad\delta((x,y))=[x+\theta].
  \]
  Irreducibility ensures that $x+\theta\ne0$. The map is a homomorphism: if a nonvertical line $y=mX+b$ meets the cubic at abscissae $x_1,x_2,x_3$, with multiplicity, evaluation of the monic intersection identity at $X=-\theta$ gives
  \[
    (x_1+\theta)(x_2+\theta)(x_3+\theta)=(b-m\theta)^2.
  \]
  The chord-and-tangent law and the vertical-line case prove the assertion. In particular, $\delta$ annihilates doubles. There is no rational two-torsion because the defining cubic is irreducible. The rational torsion group has odd order, so $\delta$ also annihilates torsion.

  The square conditions verify that $S,T$ lie on the curve. The images of $P,Q,S,T$ are $[\theta]$, $[\theta-1]$, $[\alpha]$, and $[\beta]$, independent by Lemma~\ref{lem:four-classes}. Suppose that an integral combination of these points, with coefficients not all zero, is torsion. Let $d$ be the greatest common divisor of the coefficients, and let $R$ be the combination obtained by dividing them by $d$. Since $dR$ is torsion, $R$ is torsion as well. The resulting relation is therefore primitive. Applying $\delta$ forces every coefficient to be even, contradicting primitivity. Thus the four points are independent modulo torsion.

  Now suppose that $2R=bP+cQ+eS+fT$. Applying $\delta$ makes all four coefficients even. Subtracting their half-combination from $R$ leaves rational two-torsion and therefore gives $O$. This proves saturation at two.

  For the ideals, the square conditions give $\gcd(u,6)=1$ and $\gcd(v,30)=1$. The derivatives are
  \[
    f_a'(3/4)=\frac{2u^2-137}{48},\qquad
    f_a'(15/16)=\frac{14v^2-57569}{3840},
    \quad 57569=23\cdot2503.
  \]
  The additional coprimality conditions and Lemma~\ref{lem:local-power}, with exponent two, prove the norm and square identities for $\mathfrak a,\mathfrak b$.

  The ideal identities place $[\alpha],[\beta]$ in $V_K$. If $[\mathfrak a]^e[\mathfrak b]^f=1$ for $e,f\in\{0,1\}$, Lemma~\ref{lem:virtual-units} puts $[\alpha^e\beta^f]$ in the span of $[\theta],[\theta-1]$. Lemma~\ref{lem:four-classes} forces $e=f=0$. Thus the four independent squareclasses yield two independent classes in $\Cl(K)[2]$ after quotienting by the two unit classes.
\end{proof}

\begin{samepage}
\begin{corollary}\label{cor:hilbert}
  Under the ideal-class hypotheses of Theorem~\ref{thm:four-points}, the extension
  \[
    K_a\bigl(\sqrt{\theta(\theta-1)},\sqrt{(4\theta-3)(16\theta-15)}\bigr)/K_a
  \]
  is biquadratic and unramified at every finite and infinite prime. It is contained in the ordinary Hilbert class field.
\end{corollary}
\end{samepage}

The proof implements the signature and Kummer-congruence strategy recalled in \cite[\S I.6(b--c) and Proposition I.5]{GrasGras1975}; we give the local argument for the present non-Galois fields.

\begin{proof}
  Let $h_1=\theta(\theta-1)$ and $h_2=\alpha\beta$. Both are totally positive and have even valuations at every finite prime, by the ideal identities. Hence at odd primes their quadratic extensions are unramified or split. At primes above two both are units, and
  \[
    h_1\equiv(\theta^2+1)^2\pmod{4\mathcal O_{K_a}},\qquad
    h_2\equiv1\pmod{4\mathcal O_{K_a}}.
  \]
  For a unit $h\equiv c^2\pmod4$, the polynomial $Z^2-cZ+(c^2-h)/4$ is integral and has unit discriminant $h$; its quadratic algebra is therefore unramified or split at a dyadic prime. Total positivity handles the real places. Finally, Lemma~\ref{lem:four-classes} makes $[h_1]$ and $[h_2]$ independent, so the extension has degree four.
\end{proof}

\begin{samepage}
  \begin{corollary}\label{cor:three-points}
    For the parameters in Corollary~\ref{cor:composite}, the points
    \[
      (0,1),\qquad (-1,1),\qquad (-3/4,m^{N/2}/8)
    \]
    on $E_a$ are independent modulo torsion and generate a subgroup saturated at two.
  \end{corollary}
\end{samepage}

\begin{proof}
  The norm identity verifies the third point. Lemmas~\ref{lem:units} and \ref{lem:obstruction} make their three images under $\delta$ independent. The relation and halving arguments in the proof of Theorem~\ref{thm:four-points} apply.
\end{proof}

\section{Pell parameters and a rational family}

\begin{samepage}
\begin{theorem}\label{thm:pell}
  Put $\eta=9+2\sqrt{20}$. All positive integer pairs $(u,v)$ for which \eqref{eq:two-squares} holds with an integer $a\geq3$ satisfying $a\equiv30$ or $34\pmod{52}$ are given by
  \begin{align*}
    v+u\sqrt{20} & =(79+13\sqrt{20})\eta^k,
                 & k                         & \geq0,\quad k\equiv6,7,14,15,22,27\pmod{28}, \\
    v+u\sqrt{20} & =(191+41\sqrt{20})\eta^k,
                 & k                         & \geq0,\quad k\equiv0,5,12,13,20,21\pmod{28},
  \end{align*}
  with $a=(u^2-73)/12$. There are infinitely many such parameters and infinitely many geometric isomorphism classes of the associated elliptic curves.
\end{theorem}
\end{samepage}

\begin{proof}
  For a positive solution of $v^2-20u^2=2861$, set $z=v+u\sqrt{20}$. There is a unique integer $k\geq0$ with
  \[
    \sqrt{2861}\leq z\eta^{-k}<\eta\sqrt{2861}.
  \]
  Write $z_0=z\eta^{-k}=v_0+u_0\sqrt{20}$. Both coordinates are integers because $\eta^{-1}=9-2\sqrt{20}$. The conjugate of $z_0$ is $2861/z_0$. Hence $u_0=(z_0-2861/z_0)/(2\sqrt{20})$, and the reduction interval, together with $\eta-\eta^{-1}=4\sqrt{20}$, gives
  \[
    v_0>0,\qquad 0\leq u_0<2\sqrt{2861}<107.
  \]
  Testing $v_0^2=20u_0^2+2861$ modulo $16,3,7,11$ in this proved finite interval leaves $u_0=13,19,41,47,55,79,85,97$. Modulo thirteen only $13,41,55$ remain. Their candidate squares are $79^2$, $191^2$, and $63361$, respectively; the last lies strictly between $251^2$ and $252^2$. Both valid seeds lie in the reduction interval. Thus the unrestricted nonnegative orbits of $(v_0,u_0)=(79,13),(191,41)$ give all positive solutions.

  In the coordinate order $(v,u)$, multiplication by $\eta$ gives the recurrence
  \begin{equation}\label{eq:recurrence}
    \binom{v_{k+1}}{u_{k+1}}=
    M\binom{v_k}{u_k},\qquad
    M=\begin{pmatrix}9&40\\2&9\end{pmatrix}.
  \end{equation}
  The congruences $a\equiv30\pmod{52}$ and $a\equiv34\pmod{52}$ correspond respectively to $u^2\equiv433\pmod{624}$ and $u^2\equiv481\pmod{624}$. Each seed has period $56$ under \eqref{eq:recurrence} modulo $624$; the complete calculation gives Table~\ref{tab:pell}. The matrix has determinant one, and each initial ordered pair first returns after $56$ updates. Thus the table is exhaustive. The successful exponents repeat modulo $28$, giving the stated lists. The initial permissible pair is $(v,u)=(191,41)$ and $a=134$; every subsequent permissible parameter is larger.

  \begin{table}[ht]
    \centering
    \caption{Complete congruence check for the two Pell seeds modulo $624$.}
    \label{tab:pell}
    \begin{tabular}{ccc}
      \toprule
      Seed $(v_0,u_0)$ & $a\bmod52$ & Successful exponents $0\leq k<56$ \\
      \midrule
      $(79,13)$        & $30$       & $6,15,22,27,34,43,50,55$          \\
      $(79,13)$        & $34$       & $7,14,35,42$                      \\
      $(191,41)$       & $30$       & $0,5,12,21,28,33,40,49$           \\
      $(191,41)$       & $34$       & $13,20,41,48$                     \\
      \bottomrule
    \end{tabular}
  \end{table}

  Both positive coordinates grow strictly along either orbit. Finally,
  \[
    j(E_a)=\frac{256(a^2+a+1)^3}{d_a}\sim256a^2
    \quad(a\longrightarrow\infty).
  \]
  The resulting $j$-invariants are unbounded, proving the assertion about geometric isomorphism classes.
\end{proof}

\begin{proposition}\label{prop:subsequence}
  In the orbit with seed $(v_0,u_0)=(191,41)$, take $k=28r$ with $r\geq0$. The ideal-class hypotheses of Theorem~\ref{thm:four-points} hold exactly when $r\not\equiv63\pmod{313}$. These parameters give infinitely many isomorphism classes of cubic fields.
\end{proposition}

\begin{proof}
  Theorem~\ref{thm:pell} gives the rank-four hypotheses. The remaining congruences follow from the complete periods in Table~\ref{tab:coprime}, obtained from \eqref{eq:recurrence}. In each row the initial ordered pair first returns at the displayed period; invertibility of $M$ modulo the modulus excludes a nonperiodic initial segment. The two residues in the last row say $2503\mid v_k$ exactly when $k\equiv512\pmod{1252}$. For $k=28r$ this becomes $7r\equiv128\pmod{313}$, or $r\equiv63\pmod{313}$.

  \begin{table}[ht]
    \centering
    \caption{Complete coprimality checks on the orbit from $(191,41)$.}
    \label{tab:coprime}
    \begin{tabular}{cccc}
      \toprule
      Modulus & Period & Coordinate & Zero exponents \\
      \midrule
      $137$   & $46$   & $u_k$      & none           \\
      $23$    & $8$    & $v_k$      & none           \\
      $2503$  & $2504$ & $v_k$      & $512,1764$     \\
      \bottomrule
    \end{tabular}
  \end{table}

  For completeness, the longest row has a short certificate. With $w=(191,41)^{\mathsf T}$, calculation modulo $2503$ gives
  \begin{gather*}
    M^8w=(52,993)^{\mathsf T},\qquad M^{512}w=(0,1022)^{\mathsf T},\\
    M^{1252}w=-w,\qquad M^{1764}w=(0,1481)^{\mathsf T},\qquad M^{2504}w=w.
  \end{gather*}
  Since $2504=8\cdot313$ and $313$ is prime, every proper divisor of $2504$ divides either $8$ or $1252$. The first and third identities therefore prove the period. The invariant $v^2-20u^2=2861$ has at most two states with $v=0$ modulo the prime $2503$; the displayed states give both. This verifies the zero list without a long enumeration.

  The unit equation has finitely many solutions in a fixed number field \cite[Theorem IX.4.1]{Silverman2009}. If infinitely many of these fields were isomorphic to one fixed field $K$, the images $x$ of $\theta$ would give infinitely many units for which $1-x$ is a unit. They are distinct for distinct parameters because
  \[
    a=\frac{1-x^3+x^2}{x^2-x},\qquad x\ne0,1.
  \]
  This is a contradiction, and a finite collection of fields cannot account for the infinite sequence.
\end{proof}

We now use the equation defining $E_a$ over rational function fields; the integer and positivity assumptions on $a$ in the preceding results are not imposed below.

\begin{corollary}\label{cor:rational-family}
  For an indeterminate $t$, put
  \begin{align*}
    u(t) & =\frac{41t^2-382t+820}{20-t^2},    \\
    v(t) & =\frac{191t^2-1640t+3820}{20-t^2}, \\
    a(t) & =\frac{u(t)^2-73}{12}.
  \end{align*}
  The four sections
  \[
    (0,1),\quad(-1,1),\quad(-3/4,u(t)/8),\quad(-15/16,v(t)/64)
  \]
  are independent modulo torsion in $E_{a(t)}(\Q(t))$ and generate a subgroup saturated at two. In particular, the rank over $\Q(t)$ is at least four.
\end{corollary}

\begin{proof}
  Intersecting the line $v=191-t(u+41)$ through $(u,v)=(-41,191)$ with $v^2-20u^2=2861$ gives the displayed rational parametrization at its second intersection. Direct substitution verifies the equation, and $u(0)=41$, $v(0)=191$, $a(0)=134$. The projective Weierstrass equation has coefficients in $\Q[t]_{(t)}$ and unit discriminant there, so it defines an elliptic scheme over this local ring. By properness, every $\Q(t)$-point extends over the local ring, and the group law makes specialization a homomorphism. A relation modulo torsion therefore specializes to such a relation on $E_{134}$, contrary to Theorem~\ref{thm:four-points}. This proves independence.

  There is no nonzero two-torsion over $\Q(t)$. Otherwise $f_{a(t)}$ would have a root in $\Q(t)$. Since $a(t)$ is regular at zero, this polynomial is monic over the integrally closed local ring $\Q[t]_{(t)}$. The root would belong to that ring and reduce to a rational root of $f_{134}$, contradicting irreducibility.

  If twice a rational section is an integral combination of the four sections, specialize at zero and apply the squareclass independence on $E_{134}$. All four coefficients are even. Subtracting their half-combination leaves rational two-torsion over $\Q(t)$, hence zero. This proves saturation at two.
\end{proof}

% The argument gives a rank lower bound over $\Q(t)$; the exact rank remains undetermined.

\section{An infinite rational rank-five subfamily}

To obtain a fifth point, impose $w^2=577-72a$. Together with the two earlier square conditions, this defines the curve $C\subset\mathbb P^3_{\Q}$ with homogeneous coordinates $[U:V:W:Z]$ and equations
\begin{equation}\label{eq:rank-five-base}
  V^2-20U^2=2861Z^2,\qquad W^2+6U^2=1015Z^2.
\end{equation}
On $Z\ne0$ put $u=U/Z$, $v=V/Z$, $w=W/Z$, and $a=(u^2-73)/12$. The second equation gives $w^2=577-72a$, so $(8,w)$ lies on $E_a$.

\begin{proposition}\label{prop:rank-five-rational}
  The curve $C$ is a smooth genus-one curve with infinitely many rational points. Over $\Q(C)$ the five points
  \[
    (0,1),\quad(-1,1),\quad(-3/4,u/8),\quad(-15/16,v/64),\quad(8,w)
  \]
  are independent modulo torsion and generate a subgroup saturated at two. Consequently, infinitely many geometrically nonisomorphic curves $E_a$, with $a\in\Q$, satisfy both square conditions in \eqref{eq:two-squares} and have rank at least five.
\end{proposition}

\begin{proof}
  The diagonal pencil of the two quadrics has determinant proportional to
  \[
    \lambda\mu(-20\lambda+6\mu)(-2861\lambda-1015\mu).
  \]
  Its four roots are distinct, so the Jacobian criterion shows that the intersection is smooth. As a complete intersection of two quadrics in $\mathbb P^3$, it is geometrically connected; adjunction gives $\omega_C\simeq\mathcal O_C$, hence genus one. The same calculation gives good reduction at $17,19,31$: none divides $2\cdot20\cdot6\cdot2861\cdot1015\cdot37466$.

  To show that $C(\Q)$ is infinite, we compare its rational points with the torsion bound from good reduction. Direct enumeration gives
  \[
    \#C(\F_{17})=16,\qquad \#C(\F_{19})=24,\qquad \#C(\F_{31})=28.
  \]
  These small counts can be obtained by summing $r_p(20u^2+2861)r_p(1015-6u^2)$ over $u\in\F_p$, then adding $r_p(20)r_p(-6)$ for the points at infinity, where $r_p(b)=\#\{z\in\F_p:z^2=b\}$. The affine counts are $16,24,24$, and the contributions at infinity are $0,0,4$.

  The eight choices of signs in
  \[
    [U:V:W:Z]=[\pm59:\pm1679:\pm977:31]
  \]
  give eight distinct rational points. Choose one as origin, identifying $C$ with its Jacobian. Prime-to-$p$ torsion injects under good reduction \cite[Proposition VII.3.1(b)]{Silverman2009}. Reduction at $17$ excludes every odd primary torsion subgroup except possibly the $17$-primary subgroup; reduction at $31$ excludes that subgroup as well. The two-primary torsion injects at both primes, so its order divides $\gcd(16,28)=4$. At least one of the eight rational points therefore has infinite order, and $C(\Q)$ is infinite.

  We next prove independence by specializing at $c_0=[-59:1679:977:31]$ with $a_0=-5556/961$. The five points become
  \[
    (0,1),\quad(-1,1),\quad(-3/4,-59/248),\quad
    (-15/16,1679/1984),\quad(8,977/31).
  \]
  The cubic $f_{a_0}(X)$ reduces modulo two to $X^3+X^2+1$, so it is irreducible. Let $\theta$ be a root. To certify independence of the descent values, use their squareclass representatives
  \[
    \theta,\quad\theta-1,\quad4\theta-3,\quad16\theta-15,\quad\theta+8.
  \]
  For $(p,r)=(3,2),(5,4),(7,5),(11,4),(13,2)$, the polynomial $f_{a_0}(X)$ lies in $\Z_p[X]$, and Hensel's lemma lifts its simple root $r$ modulo $p$ to a root in $\Q_p$. The resulting embedding $\Q(\theta)\hookrightarrow\Q_p$ defines a degree-one prime with residue $\theta=r$. The five displayed values have nonzero residues there. Record a square residue as $0$ and a nonsquare as $1$. With rows ordered by the pairs $(p,r)$ and columns by the five displayed elements, this gives
  \begin{equation}\label{eq:rank-five-residue-matrix}
    \begin{pmatrix}
      1 & 0 & 1 & 1 & 0 \\
      0 & 1 & 1 & 0 & 1 \\
      1 & 0 & 1 & 0 & 1 \\
      0 & 0 & 1 & 0 & 0 \\
      1 & 0 & 1 & 0 & 0
    \end{pmatrix},
    \qquad \det=1\quad\hbox{in }\F_2.
  \end{equation}
  For a direct check of invertibility, let $R_i$ denote the rows and $e_i$ the standard row vectors of $\F_2^5$. Then
  \[
    R_4=e_3,\quad R_5+R_4=e_1,\quad R_1+R_5=e_4,\quad
    R_3+R_5=e_5,\quad R_2+e_3+e_5=e_2.
  \]
  A square unit has square residue at each of these primes. The invertible matrix therefore excludes all $31$ nonempty squareclass products. The descent argument in the proof of Theorem~\ref{thm:four-points} proves independence modulo torsion of the five specialized points. This argument uses only irreducibility and nonsingularity, so it applies at the rational parameter $a_0$. The fibre at $c_0$ is nonsingular, and its smooth proper model extends specialization to every rational section. A generic relation modulo torsion would specialize to such a relation in that fibre. Thus the five sections are independent over $\Q(C)$.

  To prove saturation at two, put $F=\Q(C)$ and $\mathcal R=\mathcal O_{C,c_0}$. This local ring is an integrally closed discrete valuation ring with residue field $\Q$, and $a$ is regular at $c_0$. A root of the monic polynomial $f_a\in\mathcal R[X]$ in $F$ would belong to $\mathcal R$ and reduce to a rational root of $f_{a_0}$. Hence $E_a(F)[2]=0$. Denote the five sections by $P_1,\ldots,P_5$. If $2R=\sum_i n_iP_i$ in $E_a(F)$, specialization at $c_0$ and the invertible residue matrix force every $n_i$ to be even. The point $R-\sum_i(n_i/2)P_i$ is then two-torsion and vanishes. This proves saturation at two.

  Finally, we specialize to infinitely many fibres. The function $a$ is nonconstant on $C$, and the formula for $j(E_a)$ in the proof of Theorem~\ref{thm:pell} shows that $j$ is nonconstant as well. The specialization theorem \cite[Theorem C.20.3]{Silverman2009} therefore makes specialization injective at all but finitely many points of $C(\Q)$. The five independent sections give rank at least five on infinitely many rational fibres. The map $j:C\to\mathbb P^1$ has finite fibres, so these include infinitely many curves that are pairwise nonisomorphic over $\overline{\Q}$.
\end{proof}

\begin{remark}\label{rem:rank-five-integrality}
  The construction in Proposition~\ref{prop:rank-five-rational} has exactly one integral value of $a$. If $a\in\Z$, then $u^2=12a+73\geq0$ and $w^2=577-72a\geq0$ give $-6\leq a\leq8$. A rational number whose square is an integer is an integer. The first square condition leaves $a=-6,-4,-2,4,8$; the corresponding values of $v^2$ are $2881,3361,3841,5281,6241$, of which only $6241=79^2$ is a square. Hence $a=8$, realized by $(u,v,w)=(13,79,1)$.
\end{remark}

For integers $a$, the first square condition gives $u^2\equiv1\pmod{12}$. After changing the sign of $u$, write $u=6s+1$; then $a=3s^2+s-6$. We wonder whether there are infinitely many integers $s$ with $3s^2+s-6\geq3$ for which \(\rank E_{3s^2+s-6}(\Q)\geq5\).


\section*{Use of AI-assisted tools}

OpenAI Codex (GPT-6) was used to assist with a mathematical audit of the completed manuscript, to prepare scripts for exact algebraic checks, to check references, and to polish the exposition. The results and proofs were developed before this use of AI. The author retains responsibility for the mathematical content, references, and final text.


\bibliographystyle{alpha}
\begingroup
\raggedright
\bibliography{ennola}
\endgroup
\end{document}
